\documentclass[11pt]{article}
\usepackage[margin=1in]{geometry}
\usepackage{amsmath, amssymb, amsthm}
\usepackage{hyperref}

\newcommand{\blank}{\underline{\hspace{1.5cm}}}

\title{MAT 3150 --- Homework 3\\
\large Problem 3 Walkthrough: $\infty$ + Bounded $= \infty$}
\author{}
\date{}

\begin{document}
\maketitle

\textbf{Goal.} Prove that if $\displaystyle\lim_{n\to\infty}x_n=\infty$
and $(y_n)_n$ is bounded, then
$\displaystyle\lim_{n\to\infty}(x_n+y_n)=\infty$.

% ============================================================
\section*{The definitions}

\paragraph{Bounded.} $(y_n)$ is bounded if there is a number $M>0$ with
\[
  |y_n| \le M \quad\text{for every } n .
\]
Unpacking the absolute value: $-M \le y_n \le M$. For this problem you
only need the \emph{lower} half, $y_n \ge -M$. (Why the lower half? We
want $x_n+y_n$ to be \emph{big}, so the danger is $y_n$ being very
negative.)

\paragraph{Diverging to $\infty$.} $a_n\to\infty$ means:
\begin{quote}
For every $K\in\mathbb{R}$ there is $N\in\mathbb{N}$ such that
$a_n > K$ for all $n\ge N$.
\end{quote}
Same game as the $\varepsilon$ definition, but now the challenger hands
you a \emph{height} $K$ (maybe huge, like $10^{6}$), and you must find a
point $N$ after which every term is above $K$.

\paragraph{The key move for this whole topic.} The hypothesis
$x_n\to\infty$ works for \emph{every} height. So you are free to apply it
to a \emph{different} height than $K$ --- whatever height makes your
target work out.

% ============================================================
\section*{Worked example: if $x_n\to\infty$, then $3x_n-5\to\infty$}

\paragraph{Scratch work.} Given a height $K$, we want $3x_n-5>K$. Solve
for $x_n$:
\[
  3x_n - 5 > K \iff x_n > \frac{K+5}{3}.
\]
So we should apply the hypothesis to the height $\dfrac{K+5}{3}$, not to
$K$.

\paragraph{The proof.}
\begin{proof}
Let $K\in\mathbb{R}$. Since $x_n\to\infty$, applying the definition with
the height $\dfrac{K+5}{3}$ gives $N\in\mathbb{N}$ such that
\[
  x_n > \frac{K+5}{3} \quad\text{for all } n\ge N .
\]
Then for all $n\ge N$,
\[
  3x_n - 5 > 3\cdot\frac{K+5}{3} - 5 = K .
\]
Since $K$ was arbitrary, $3x_n-5\to\infty$.
\end{proof}

\noindent\textbf{Notice:} the proof starts with ``Let $K$'' (the height
for the \emph{conclusion}), then feeds a \emph{modified} height into the
\emph{hypothesis}.

% ============================================================
\newpage
\section*{Your turn: $x_n\to\infty$ and $(y_n)$ bounded $\Rightarrow$ $x_n+y_n\to\infty$}

\paragraph{1. Write down what you're given.}
\begin{itemize}
  \item Bounded: there is $M>0$ with $\blank$ for every $n$. \\
        In particular, $y_n \ge \blank$ for every $n$.
  \item $x_n\to\infty$: for every height $H$ there is $N$ with
        $\blank$ for all $n\ge N$.
\end{itemize}

\paragraph{2. Write down what you must show.}
For every $K$ there is $N$ such that $\blank > K$ for all $n\ge N$.

\paragraph{3. Scratch work: get rid of $y_n$.}
Using $y_n\ge -M$, find a lower bound for $x_n+y_n$ that involves only
$x_n$ and $M$:
\[
  x_n + y_n \;\ge\; \blank .
\]
For this lower bound to be bigger than $K$, you need
\[
  x_n > \blank .
\]
That is the height to feed into the hypothesis.

\vspace{1cm}

\paragraph{4. The proof.} Fill in the template:
\begin{quote}
Since $(y_n)$ is bounded, there is $M>0$ with $|y_n|\le M$, so in
particular $y_n\ge -M$ for all $n$.

Let $K\in\mathbb{R}$. Since $x_n\to\infty$, applying the definition with
the height $\blank$ gives $N\in\mathbb{N}$ such that
\[
  x_n > \blank \quad\text{for all } n\ge N .
\]
Then for all $n\ge N$,
\[
  x_n + y_n \;\ge\; \blank \;>\; \blank \;=\; K .
\]
Since $K$ was arbitrary, $x_n+y_n\to\infty$.
\end{quote}

\paragraph{Check yourself.}
\begin{itemize}
  \item Why is the first step ``$\ge$'' and not ``$>$''?
  \item Which step uses boundedness, and which uses $x_n\to\infty$?
  \item Where did you use that $K$ was arbitrary?
\end{itemize}

\end{document}
